%%
%% This is file `sample-acmtog.tex',
%% generated with the docstrip utility.
%%
%% The original source files were:
%%
%% samples.dtx  (with options: `all,journal,bibtex,acmtog')
%%
%% IMPORTANT NOTICE:
%%
%% For the copyright see the source file.
%%
%% Any modified versions of this file must be renamed
%% with new filenames distinct from sample-acmtog.tex.
%%
%% For distribution of the original source see the terms
%% for copying and modification in the file samples.dtx.
%%
%% This generated file may be distributed as long as the
%% original source files, as listed above, are part of the
%% same distribution. (The sources need not necessarily be
%% in the same archive or directory.)
%%
%%
%% Commands for TeXCount
%TC:macro \cite [option:text,text]
%TC:macro \citep [option:text,text]
%TC:macro \citet [option:text,text]
%TC:envir table 0 1
%TC:envir table* 0 1
%TC:envir tabular [ignore] word
%TC:envir displaymath 0 word
%TC:envir math 0 word
%TC:envir comment 0 0
%%
%% The first command in your LaTeX source must be the \documentclass
%% command.
%%
%% For submission and review of your manuscript please change the
%% command to \documentclass[manuscript, screen, review]{acmart}.
%%
%% When submitting camera ready or to TAPS, please change the command
%% to \documentclass[sigconf]{acmart} or whichever template is required
%% for your publication.
%%
%%
\documentclass[acmtog,nonacm]{acmart}
\usepackage{multirow}
\usepackage{placeins}
\usepackage{cleveref}
\usepackage{indentfirst}
%%
%% \BibTeX command to typeset BibTeX logo in the docs
\AtBeginDocument{%
  \providecommand\BibTeX{{%
    Bib\TeX}}}

%% Rights management information.  This information is sent to you
%% when you complete the rights form.  These commands have SAMPLE
%% values in them; it is your responsibility as an author to replace
%% the commands and values with those provided to you when you
%% complete the rights form.
\setcopyright{none}
\copyrightyear{2018}
\acmYear{2018}
\acmDOI{}

%%
%% These commands are for a JOURNAL article.
\acmJournal{TOG}
\acmVolume{37}
\acmNumber{4}
\acmMonth{8}

%%
%% Submission ID.
%% Use this when submitting an article to a sponsored event. You'll
%% receive a unique submission ID from the organizers
%% of the event, and this ID should be used as the parameter to this command.
%%\acmSubmissionID{123-A56-BU3}

%%
%% For managing citations, it is recommended to use bibliography
%% files in BibTeX format.
%%
%% You can then either use BibTeX with the ACM-Reference-Format style,
%% or BibLaTeX with the acmnumeric or acmauthoryear sytles, that include
%% support for advanced citation of software artefact from the
%% biblatex-software package, also separately available on CTAN.
%%
%% Look at the sample-*-biblatex.tex files for templates showcasing
%% the biblatex styles.
%%

%%
%% The majority of ACM publications use numbered citations and
%% references.  The command \citestyle{authoryear} switches to the
%% "author year" style.
%%
%% If you are preparing content for an event
%% sponsored by ACM SIGGRAPH, you must use the "author year" style of
%% citations and references.
\citestyle{acmauthoryear}


%%
%% end of the preamble, start of the body of the document source.
\begin{document}

%%
%% The "title" command has an optional parameter,
%% allowing the author to define a "short title" to be used in page headers.
\title{Recognizing Portuguese Varieties and Rewriting Between Them with Fine-tuned Language Models}

%%
%% The "author" command and its associated commands are used to define
%% the authors and their affiliations.
%% Of note is the shared affiliation of the first two authors, and the
%% "authornote" and "authornotemark" commands
%% used to denote shared contribution to the research.
\author{\texorpdfstring{\NoCaseChange{Rodrigo Laia}}{Rodrigo Laia}}
%\authornote{Both authors contributed equally to this research.}
\email{rodrigo.laia@tecnico.ulisboa.pt}
%\orcid{1234-5678-9012}
%\author{G.K.M. Tobin}
%\email{webmaster@marysville-ohio.com}
\affiliation{%
  \institution{Instituto Superior Técnico}
  \city{Lisbon}
  %\state{Ohio}
  \country{Portugal}
}

% \author{Lars Th{\o}rv{\"a}ld}
% \affiliation{%
%   \institution{The Th{\o}rv{\"a}ld Group}
%   \city{Hekla}
%   \country{Iceland}}
% \email{larst@affiliation.org}

%%
%% By default, the full list of authors will be used in the page
%% headers. Often, this list is too long, and will overlap
%% other information printed in the page headers. This command allows
%% the author to define a more concise list
%% of authors' names for this purpose.
\renewcommand{\shortauthors}{Laia}

%%
%% The abstract is a short summary of the work to be presented in the
%% article.
\begin{abstract}
Portuguese has standardized varieties such as European Portuguese (EP) and Brazilian Portuguese (BP), yet web-scale Portuguese is dominated by BP, which can make generic Portuguese systems less reliable for EP. We study how pretrained encoder-decoder models can be adapted for both Portuguese variety identification and bidirectional EP/BP rewriting, comparing encoder-controlled and decoder-controlled formulations within a staged pipeline of supervised adaptation and reward-based continuation. Across the supervised comparisons, supervision composition produces the largest rewriting changes. GPT-Wikipedia, constructed around localized variety adaptation, gives the strongest Golden Collection rewriting results, whereas adding the Few-shot Region-aware Machine Translation (FRMT) training split sharply lowers Golden Collection performance in both directions and changes FRMT rewriting in a direction- and control-dependent manner. Control placement also depends on rewriting direction, with decoder control stronger in every BP-to-EP comparison and EP-to-BP results remaining mixed. Reward-based continuation has modest effects after GPT-Wikipedia specialization but recovers a substantial part of the Golden Collection loss after GPT-Wikipedia + FRMT specialization. These results show that Portuguese variety adaptation depends strongly on whether supervision reflects localized EP/BP rewriting and on evaluation signals that distinguish necessary changes, appropriate copying, and broader paraphrase.
\end{abstract}

%%
%% The code below is generated by the tool at http://dl.acm.org/ccs.cfm.
%% Please copy and paste the code instead of the example below.
%%
% \begin{CCSXML}
% <ccs2012>
%  <concept>
%   <concept_id>00000000.0000000.0000000</concept_id>
%   <concept_desc>Do Not Use This Code, Generate the Correct Terms for Your Paper</concept_desc>
%   <concept_significance>500</concept_significance>
%  </concept>
%  <concept>
%   <concept_id>00000000.00000000.00000000</concept_id>
%   <concept_desc>Do Not Use This Code, Generate the Correct Terms for Your Paper</concept_desc>
%   <concept_significance>300</concept_significance>
%  </concept>
%  <concept>
%   <concept_id>00000000.00000000.00000000</concept_id>
%   <concept_desc>Do Not Use This Code, Generate the Correct Terms for Your Paper</concept_desc>
%   <concept_significance>100</concept_significance>
%  </concept>
%  <concept>
%   <concept_id>00000000.00000000.00000000</concept_id>
%   <concept_desc>Do Not Use This Code, Generate the Correct Terms for Your Paper</concept_desc>
%   <concept_significance>100</concept_significance>
%  </concept>
% </ccs2012>
% \end{CCSXML}

% \ccsdesc[500]{Do Not Use This Code~Generate the Correct Terms for Your Paper}
% \ccsdesc[300]{Do Not Use This Code~Generate the Correct Terms for Your Paper}
% \ccsdesc{Do Not Use This Code~Generate the Correct Terms for Your Paper}
% \ccsdesc[100]{Do Not Use This Code~Generate the Correct Terms for Your Paper}

%%
%% Keywords. The author(s) should pick words that accurately describe
%% the work being presented. Separate the keywords with commas.
\keywords{Portuguese varieties, variety identification, text rewriting, encoder-decoder models, parameter-efficient fine-tuning}

% \received{20 February 2007}
% \received[revised]{12 March 2009}
% \received[accepted]{5 June 2009}

%%
%% This command processes the author and affiliation and title
%% information and builds the first part of the formatted document.
\maketitle

\section{Introduction}
Most Natural Language Processing (NLP) systems expose Portuguese as a single
language option even though it is used through different national standards.
When web-scale training data underrepresent one standard, variety-sensitive
behavior can drift toward the variety that is more frequent in the data. BP is
substantially more visible online than EP, which can make generic Portuguese
systems default toward BP and serve EP less reliably
\citep{Sousa_Almeida_Silvano_Cantante_Campos_Jorge_2025}.

Portuguese is shared across communities with different linguistic standards,
and the EP/BP contrast involves orthographic conventions, lexical choices,
morphosyntactic preferences, clitic placement, verbal periphrases, and forms of
address \citep{barreiro-mota-2018-paraphrastic}. These distinctions provide
evidence for variety identification and determine which changes are appropriate
during rewriting, where many valid EP/BP adaptations require only localized
edits and some sentences require no modification.

Variety identification tests whether a system can detect subtle variety cues,
while bidirectional EP/BP rewriting tests whether it can produce the requested
variety while preserving the original meaning. Rewriting makes supervision
quality especially important, as aligned EP/BP data must distinguish
variety-specific adaptation from paraphrastic or editorial changes. The
available resources differ in domain, alignment quality, construction procedure,
and the transformations represented by their aligned pairs. Fine-tuning a shared
model on these resources therefore exposes it to different rewriting behavior
depending on how the supervision is constructed and combined.

The experiments fine-tune pretrained T5Gemma 2 encoder-decoder models for both
tasks under a shared control-string setup. Encoder-controlled and
decoder-controlled formulations differ in whether variety information is placed
in the encoder input or the decoder target. Training follows a staged pipeline
with large-scale supervised adaptation, targeted supervised specialization, and
reward-based continuation using a group-relative objective inspired by Group
Relative Policy Optimization (GRPO)
\citep{shao2024deepseekmathpushinglimitsmathematical}. A lower-cost supervised
protocol-selection study fully fine-tunes the 270M-parameter model to define the
training protocol used in the subsequent experiments, which adapt the
4B-parameter model with LoRA.

The paper makes three contributions.
\begin{itemize}
\item It evaluates a shared T5Gemma 2 setup for EP/BP variety identification and bidirectional rewriting under encoder-controlled and decoder-controlled formulations.
\item It constructs and uses GPT-Wikipedia as reviewed supervision designed around localized EP/BP adaptation and examines how supervision composition changes identification and rewriting.
\item It studies staged adaptation from large-scale OpenSubtitles supervision through targeted specialization and reward-based continuation, combining aggregate evaluation with analysis of copying, required edits, and unnecessary reformulation.
\end{itemize}

Supervision composition produces the largest rewriting changes: GPT-Wikipedia
aligns most closely with the conservative Golden Collection, while adding FRMT
sharply lowers Golden Collection rewriting and changes FRMT performance
differently across directions and control formulations. Control placement also depends on rewriting direction, with decoder control
stronger in every reported BP-to-EP rewriting comparison, whereas EP-to-BP
results are mixed and include settings where encoder control performs better. Stage A and Stage C effects
depend on the surrounding supervision, with reward-based continuation producing
its clearest Golden Collection recovery after GPT-Wikipedia + FRMT
specialization. Copy analysis further shows that automatic metrics can reward
appropriate preservation but cannot fully distinguish required adaptation from
acceptable paraphrase or unnecessary reformulation.

The remainder of the paper is organized as follows. \Cref{sec:related-work} reviews the most relevant literature. \Cref{sec:methods} describes the data, task formulations, and training setup. \Cref{sec:results} presents the quantitative and qualitative results. \Cref{sec:discussion-limitations} discusses the main trade-offs and limitations, and \Cref{sec:conclusion} concludes the paper.

\section{Related Work}
\label{sec:related-work}

This section aims to position this research within prior work on Portuguese
variety processing and model adaptation. \Cref{sec:rw-portuguese-variety}
reviews work on European and Brazilian Portuguese identification and
rewriting, while \Cref{sec:rw-adaptation-controlled-generation} examines
shared encoder-decoder modeling, staged adaptation, and reward-based
continuation.

\subsection{Portuguese Variety Identification and Rewriting}
\label{sec:rw-portuguese-variety}

Early work on EP/BP identification showed that orthographic, lexical, and
character-level patterns provide strong discriminative evidence. Using
character \(n\)-grams and word \(n\)-gram language models,
\citet{DBLP:conf/konvens/ZampieriG12} demonstrated the feasibility of
distinguishing the two varieties from surface cues. \citet{preda-etal-2024-across} examine whether this evidence remains reliable
across different input lengths and domains by comparing surface-based
classifiers with LoRA-adapted models. Their LoRA-adapted Albertina classifiers
obtain the strongest in-domain results, whereas a logistic-regression bigram
model performs best in the joint cross-dataset evaluation, showing that the
strongest in-domain system is not also the strongest under their cross-domain
comparison.

Cross-domain robustness also depends on what signals the training labels
encourage a model to learn. Labels inferred from country, domain, or URL
metadata can encode source-specific proxies rather than reliable variety cues
\citep{zampieri-etal-2014-report}. PtBrVarId, introduced by
\citet{Sousa_Almeida_Silvano_Cantante_Campos_Jorge_2025}, evaluates EP/BP
identification explicitly across domains. Its experiments show that partial
delexicalization can improve identification by reducing reliance on named
entities and domain-specific lexical signals, with its strongest systems
combining this procedure with BERTimbau \citep{souza2020bertimbau}. Sousa
et al. therefore provide the closest dedicated EP/BP identification reference,
whereas the present setting learns identification in the same model used for
rewriting. Their work establishes a strong dedicated identification comparison
but leaves open how identification behaves when learned alongside sentence
rewriting.

Rewriting introduces the additional requirement of preserving the source
meaning while making only the changes appropriate to the requested variety.
Early BP-to-EP work treated adaptation as a dedicated conversion problem
\citep{marujo-etal-2011-bp2ep}, while linguistic analyses show that localized
EP/BP differences can extend beyond individual words to multiword,
phrase-level, and morphosyntactic alternatives
\citep{barreiro-mota-2018-paraphrastic}. This range of possible changes
supports a sequence-to-sequence formulation without implying that every input
requires structural modification. \citet{costa-jussa-etal-2018-neural} show
that a neural sequence-to-sequence system improves over a phrase-based approach
in both Portuguese-variety directions and that human evaluators favor the
neural outputs. More recently,
\citet{sanches2024brazilianportugueseeuropeanportuguese} provide the closest
dedicated rewriting comparison by evaluating pretrained BP-to-EP systems on the
manually curated Golden Collection. Their experiments show that modern models
can perform competitively on the task, while the benchmark's conservative
references make unchanged-source copying an important reference for judging
whether a system rewrites only when needed. Portuguese identification and
rewriting therefore have strong separate foundations, while less is known
about how the two tasks behave when learned in the same model under different
forms of variety-specific supervision.

\subsection{Language Model Adaptation for Controlled Generation}
\label{sec:rw-adaptation-controlled-generation}

Outside Portuguese, \citet{khered-etal-2025-multi} provide a methodological
precedent by combining Arabic dialect identification and
dialect-to-Modern-Standard-Arabic translation within a multi-task
encoder-decoder model. Their experiments show that the shared setup, together
with the additional training data used in the study, improves translation over
the evaluated baselines. The reported gain concerns this overall design and
therefore provides evidence for the combined setup rather than an isolated
multi-task effect.

Sharing the tasks addresses the model interface, while a separate question is
how supervision of different scale and quality should be introduced during
adaptation. Work on staged Machine Translation (MT) distinguishes broad
exposure from smaller amounts of high-quality parallel data or targeted
supervised specialization
\citep{xu2024paradigmshiftmachinetranslation,alves2024tower}.
The Tower+ and Omnilingual MT models extend this progression with preference,
reward-based, or other post-training phases
\citep{rei-etal-2026-tower,omnilingualmtteam2026omnilingualmtmachinetranslation}.
Across these systems, broad adaptation, targeted supervision, and post-training
serve distinct functions, supporting their separation when the available data
differ in scale and quality.

Reward-based updates have also been used to refine already supervised
translation models in scalar-feedback and low-resource dialect settings
\citep{nguyen-etal-2017-reinforcement,thu2023rlnmt}, establishing that
task-specific feedback can guide post-training continuation.

Task-specific feedback can be defined through reference-free quality estimates,
learned translation preferences, or fine-grained error information, each of
which can provide task-relevant MT rewards
\citep{he-etal-2024-improving,xu2024advancingtranslationpreferencemodeling,ramos-etal-2026-fine}.
However, reference-free reward optimization can become misaligned with actual
translation quality \citep{he-etal-2024-improving}, making reward design
central. Recent systems therefore combine or decompose rewards across
complementary criteria \citep{feng-etal-2025-mt-r1,10.1162/TACL.a.65}. For EP/BP rewriting, the reward signal must favor required variety-specific changes without encouraging broader reformulation, while preserving the source meaning.

\section{Experimental Setup}
\label{sec:methods}

This section aims to define the experimental setup used to study Portuguese
variety identification and rewriting. \Cref{sec:data-task-construction} introduces the data resources and explains
how aligned EP/BP pairs are converted into task-specific rewriting and
identification examples, \Cref{sec:task-control-formulations} describes the
model scales and control formulations, and \Cref{sec:training-procedure}
presents the staged training procedure and optimization settings. \Cref{sec:method-protocol-selection}
defines the supervised protocol-selection study, while
\Cref{sec:evaluation-setup-paper} describes the evaluation procedure.

\subsection{Data and Task Construction}
\label{sec:data-task-construction}

The experiments draw on four aligned EP/BP resources whose different
characteristics support broad adaptation, targeted specialization, and
evaluation. OpenSubtitles \citep{4992de1b5fb34f3e9691772606b36edf}
supplies large-scale exposure through aligned movie and television dialogue.
Its scale supports broad EP/BP adaptation, although its pairs are noisier and
less controlled than the targeted supervision resources.

More controlled supervision comes from GPT-Wikipedia, a resource constructed
as part of this work from Wikipedia-inspired encyclopedic contexts. Its near-literal EP/BP pairs are designed to favor localized variety
changes rather than broad paraphrase. During construction, differences not
required for variety adaptation were made identical where appropriate,
retaining equal pairs that provide copying supervision when no change is
needed.

FRMT adds region-aware variation through a machine-translation resource and
benchmark derived from English source sentences with regional outputs,
including EP and BP \citep{riley-etal-2023-frmt}. FRMT is used primarily as a
held-out evaluation resource, while its training split is added only in the
GPT-Wikipedia + FRMT condition to examine how this additional supervision
changes identification and rewriting performance. The Golden Collection
complements FRMT with a manually curated rewriting benchmark used only for
evaluation \citep{sanches2024brazilianportugueseeuropeanportuguese}. Many of its
references are identical or close to the source, so the benchmark rewards
conservative, localized rewriting.

Their relative scales and split roles are summarized before task construction
in \Cref{tab:resource-pair-counts}.

\begin{table}[H]
  \centering
  \small
  \renewcommand{\arraystretch}{1.12}
  \setlength{\tabcolsep}{4pt}
  \caption{Aligned EP/BP pair counts by resource and split before task-specific rendering.}
  \label{tab:resource-pair-counts}
  \begin{tabular}{@{}lccc@{}}
    \toprule
    \textbf{Resource} & \textbf{Train} & \textbf{Validation} & \textbf{Test} \\
    \midrule
    OpenSubtitles & 10,347,883 & 0 & 0 \\
    GPT-Wikipedia & 4,519 & 564 & 566 \\
    FRMT & 2,502 & 20 & 2,611 \\
    Golden Collection & 0 & 0 & 500 \\
    \bottomrule
  \end{tabular}
\end{table}

The OpenSubtitles counts describe the processed base resource, which has no
separate validation split. A 200-instance validation view is sampled from its
processed training data when the Stage A data are prepared, whereas the FRMT
training split supplies the additional Stage B supervision and its held-out
test split is used for evaluation. These resource-level counts precede task
construction. Each aligned pair yields one BP-to-EP and one EP-to-BP rewriting
example, while a non-equal pair also yields one identification example for each
variety. Equal pairs remain useful for rewriting and are excluded from
identification, since identical text cannot provide unambiguous supervision
for two opposing variety labels.
Additional details on resource construction and filtering are provided in
\Cref{sec:appendix-data-construction}.

\subsection{Models and Control Formulations}
\label{sec:task-control-formulations}

The experiments use T5Gemma 2 encoder-decoder models with 270M and 4B
parameters \citep{zhang2025t5gemma2seeingreading}. Their text-to-text
architecture provides a single interface for short variety-label outputs and
full-sentence rewriting. The 270M-parameter model is fully fine-tuned in a
lower-cost supervised protocol-selection study that defines the training
protocol used in the subsequent experiments. Those experiments adapt the
4B-parameter model with LoRA \citep{hu2022lora}, using rank 48,
\(\alpha=96\), and dropout 0.05, which yields 187,772,928 trainable adapter
parameters while the base parameters remain frozen.

Across both tasks, \texttt{PT} labels an EP source sentence and \texttt{BR}
labels a BP source sentence. The two control formulations differ in where
this input-variety information is represented. In encoder-controlled
rewriting, the relevant label prefixes the encoder input and the decoder target
contains only the rewritten sentence. Encoder-controlled identification
instead introduces \texttt{CLS} before the input sentence, while the decoder
predicts \texttt{PT} or \texttt{BR}.

The decoder-controlled formulation moves the input-variety label to the
generated sequence. Its rewriting inputs contain only the source sentence,
while each target begins with \texttt{PT} or \texttt{BR} and continues with
the rewrite. For identification, the same unprefixed input is paired with a
label-only target. For decoder-controlled identification, loss is computed only
on the \texttt{PT} or \texttt{BR} label token, whereas rewriting uses the
complete target sequence. \Cref{sec:method-protocol-selection} specifies the
supervised training protocol used with these task formulations, including task
mixing, tokenizer treatment, and classification scoring.

\begin{table*}[t]
  \centering
  \small
  \renewcommand{\arraystretch}{1.18}
  \setlength{\tabcolsep}{5pt}
  \caption{Sequence-to-sequence task templates. \texttt{E} and \texttt{D}
  denote encoder-controlled and decoder-controlled formulations,
  respectively. \(x_{\mathrm{PT}}\) and \(x_{\mathrm{BR}}\) are aligned EP
  and BP sentences, while \(x\) is an input sentence whose variety is to be
  identified.}
  \label{tab:training-templates}
  \begin{tabular}{@{}lccc@{}}
    \toprule
    \textbf{Setup} & \textbf{Direction} & \textbf{Encoder input} &
    \textbf{Decoder target} \\
    \midrule
    \multirow{2}{*}{\texttt{E} rewriting} &
    BP-to-EP & \texttt{BR} \(x_{\mathrm{BR}}\) &
    \(x_{\mathrm{PT}}\) \\
    & EP-to-BP & \texttt{PT} \(x_{\mathrm{PT}}\) &
    \(x_{\mathrm{BR}}\) \\
    \addlinespace[0.35em]
    \texttt{E} identification &
    --- & \texttt{CLS} \(x\) &
    \texttt{PT} or \texttt{BR} \\
    \midrule
    \multirow{2}{*}{\texttt{D} rewriting} &
    BP-to-EP & \(x_{\mathrm{BR}}\) &
    \texttt{BR} \(x_{\mathrm{PT}}\) \\
    & EP-to-BP & \(x_{\mathrm{PT}}\) &
    \texttt{PT} \(x_{\mathrm{BR}}\) \\
    \addlinespace[0.35em]
    \texttt{D} identification &
    --- & \(x\) &
    \texttt{PT} or \texttt{BR} \\
    \bottomrule
  \end{tabular}
\end{table*}

\subsection{Staged Training}
\label{sec:training-procedure}

Supervised adaptation uses the identification and rewriting examples described
above, progressing from large-scale OpenSubtitles exposure in Stage A to
targeted specialization in Stage B. In Stage A, OpenSubtitles exposes the
models to large-scale aligned EP/BP supervision. After task construction, the
Stage A pool contains 20,695,766 rewriting examples and 7,666,883
identification examples. The fixed optimization budgets expose the models to
only a small fraction of this pool, so Stage A examples do not repeat before the
maximum step budget is reached. The contribution of this warm-up is isolated
through matched configurations that omit Stage A, start directly at Stage B,
and hold the subsequent specialization setup fixed.

Stage B moves from large-scale subtitle exposure to targeted supervised
specialization. GPT-Wikipedia provides the primary targeted condition, while
the GPT-Wikipedia + FRMT condition adds the complete FRMT training split to
examine how additional region-aware supervision changes identification and
rewriting. The two training splits are concatenated without weighting or
resampling, while the task templates and training procedure remain unchanged.
After task construction, the GPT-Wikipedia condition contains 9,038 rewriting
examples and 6,652 identification examples. The GPT-Wikipedia + FRMT condition
contains 14,042 rewriting examples and 11,616 identification examples. For the
4B-parameter model, both Stage B training sets are smaller than the maximum
training-example presentation budget and can therefore repeat across shuffled
epochs. For the 270M-parameter model, the corresponding budget exceeds the
GPT-Wikipedia set but not the GPT-Wikipedia + FRMT set.

Stage C preserves the two Stage B data conditions, GPT-Wikipedia and
GPT-Wikipedia + FRMT, but continues each Stage B checkpoint using a more
conservative rewriting-only subset. The subsets are selected to favor
localized EP/BP changes, structurally close pairs, and unchanged examples where
copying is appropriate. The GPT-Wikipedia Stage C subset contains 6,686
rewriting examples balanced across directions. The GPT-Wikipedia + FRMT Stage C
subset contains 6,828 balanced rewriting examples, comprising the same 6,686
GPT-Wikipedia examples and 142 retained FRMT examples. This combined set
underlies the Stage C rows for the GPT-Wikipedia + FRMT condition in the result
tables.

For each Stage C input, four candidate rewrites are sampled and the reference
rewrite is added to the group. Candidate rewards are compared within this
group to form the policy update. For each candidate, the BLEU \citep{papineni2002bleu} and TER
\citep{snover2006study} components compare its agreement with the reference
against the score obtained by returning the source unchanged, making the reward
relative to a copy baseline rather than to the reference score alone. Under
encoder control, the reward assigns equal weight to the relative BLEU and TER
components, \(0.5R_{\mathrm{BLEU}}+0.5R_{\mathrm{TER}}\). Under decoder
control, it uses
\(0.4R_{\mathrm{BLEU}}+0.4R_{\mathrm{TER}}+0.2R_{\mathrm{label}}\), where
\(R_{\mathrm{label}}\) rewards a correct first \texttt{PT}/\texttt{BR} label.
After removing any decoder control label and normalizing whitespace, 0.10 is
subtracted when a candidate exactly copies the source even though the reference
requires a change. Since Stage C
combines several reward components, supplementary 270M reward-design
diagnostics examine the effects of individual components and candidate-grouping
choices.
\Cref{sec:appendix-rl-diagnostics} reports their design and results.

The scale- and stage-dependent optimization settings are presented in
\Cref{tab:training-settings}.

\begin{table}[H]
  \centering
  \small
  \renewcommand{\arraystretch}{1.12}
  \setlength{\tabcolsep}{2.2pt}

  \caption{Central optimization settings by stage and model. FT denotes full fine-tuning.}
  \label{tab:training-settings}

  \begin{tabular}{@{}clccc@{}}
    \toprule
    \textbf{Stage}
    & \shortstack[l]{\textbf{Model /}\\\textbf{adaptation}}
    & \shortstack{\textbf{Learning}\\\textbf{rate}}
    & \shortstack{\textbf{Warm-up}\\\textbf{steps}}
    & \shortstack{\textbf{Maximum}\\\textbf{steps}} \\
    \midrule
    A & 270M full FT & \(2 \times 10^{-5}\) & 200 & 2,500 \\
    A & 4B LoRA     & \(1 \times 10^{-5}\) & 400 & 2,500 \\
    B & 270M full FT & \(2 \times 10^{-5}\) & 150 & 1,200 \\
    B & 4B LoRA     & \(1 \times 10^{-5}\) & 150 & 1,200 \\
    C & 4B LoRA     & \(5 \times 10^{-6}\) & 20  & 600 \\
    \bottomrule
  \end{tabular}
\end{table}

All five settings use AdamW \citep{Loshchilov2017DecoupledWD} with a linear learning-rate schedule, with weight
decay set to 0.01 in supervised Stages A and B and to 0 in Stage C. Supervised
Stages A and B use validation-based early stopping, with the
lowest-validation-loss checkpoint restored.

\subsection{Supervised Protocol Selection}
\label{sec:method-protocol-selection}

Four supervised protocols are evaluated with the fully fine-tuned
270M-parameter decoder-controlled model. All four variants follow the same
OpenSubtitles Stage A and GPT-Wikipedia Stage B path, with the same T5Gemma 2
checkpoint, task construction, and comparable step budgets, and are evaluated
with the common procedure described in \Cref{sec:evaluation-setup-paper}. Only
the supervised protocol component changes, and the selected protocol is then
used in the subsequent experiments with the 4B-parameter model.

The base protocol combines rewriting and identification rows in the same
shuffled training mixture, so their contribution follows their natural
frequency rather than an enforced task balance. \texttt{PT} and \texttt{BR}
use their existing base-vocabulary representations rather than newly registered
tokenizer tokens. For identification examples, loss is computed at the
supervised label position against the normal decoder vocabulary, while
rewriting uses the complete target sequence and normal decoder vocabulary. The
task-balanced variant changes only batch composition by enforcing equal
rewriting and identification representation in each optimizer update. The
added-token variant keeps the same control strings but explicitly registers
\texttt{PT} and \texttt{BR} with the tokenizer. The
restricted-classification-loss variant changes only identification scoring by
restricting the competing output labels to \texttt{PT} and \texttt{BR}, while
rewriting remains unchanged. The results of this comparison are presented in
\Cref{sec:supervised-protocol-selection}.

\subsection{Evaluation}
\label{sec:evaluation-setup-paper}

The two tasks are evaluated under both control formulations on FRMT and the
Golden Collection. Accuracy and macro-F1 measure identification, while BLEU
\citep{papineni2002bleu}, computed with SacreBLEU \citep{post-2018-call}, and
TER \citep{snover2006study} measure rewriting. The principal aggregate
rewriting comparisons use corpus BLEU and TER, reported separately for BP-to-EP
and EP-to-BP, whereas sentence-level BLEU and TER are reserved for diagnostic
analyses and also contribute to the Stage C reward described in
\Cref{sec:training-procedure}. Rewriting uses beam search with beam size 4.
Before the rewriting scores are computed, control labels are removed from
predictions and references so that the metrics evaluate sentence content
rather than the control tokens. The copy baseline returns the source sentence
unchanged and provides an informative reference for rewriting, since some EP/BP
examples require no modification, with \Cref{sec:qualitative-analysis}
examining this behavior alongside representative outputs.

\section{Results}
\label{sec:results}

This section aims to determine how supervised protocol choice, control
placement, training-data composition, staged adaptation, and reward-based
continuation affect identification and rewriting.
\Cref{sec:supervised-protocol-selection} reports the protocol-selection
results, \Cref{sec:main-supervised-results} compares supervised identification
and rewriting after Stage B, and \Cref{sec:reward-continuation-results}
examines reward-based continuation. \Cref{sec:qualitative-analysis} analyzes
representative outputs and copying behavior.

\subsection{Supervised Protocol Selection}
\label{sec:supervised-protocol-selection}

The 270M-parameter protocol-selection study evaluates which supervised protocol
provides the most suitable basis for the subsequent experiments with the
4B-parameter model. \Cref{tab:protocol-ablation-270m-paper}
compares the four candidates after the same OpenSubtitles Stage A and
GPT-Wikipedia Stage B path.

\begin{table*}[tbp]
\centering
\small
\renewcommand{\arraystretch}{1.08}
\setlength{\tabcolsep}{3pt}
\caption{Decoder-controlled 270M full-fine-tuning supervised protocol-selection results. Macro-F1 is reported as a percentage, BLEU on a 0--100 scale, and lower TER is better.}
\label{tab:protocol-ablation-270m-paper}
\begin{tabular}{p{3.4cm}cccccccccc}
\toprule
\multirow{3}{*}{\textbf{Protocol variant}} & \multicolumn{5}{c}{\textbf{FRMT}} & \multicolumn{5}{c}{\textbf{Golden Collection}} \\
\cmidrule(lr){2-6} \cmidrule(lr){7-11}
& \multirow{2}{*}{\shortstack{\textbf{Macro-F1}\\\textbf{(\%)}}} & \multicolumn{2}{c}{\shortstack{\textbf{BLEU}\\\textbf{(0--100)}}} & \multicolumn{2}{c}{\textbf{TER}} & \multirow{2}{*}{\shortstack{\textbf{Macro-F1}\\\textbf{(\%)}}} & \multicolumn{2}{c}{\shortstack{\textbf{BLEU}\\\textbf{(0--100)}}} & \multicolumn{2}{c}{\textbf{TER}} \\
\cmidrule(lr){3-4} \cmidrule(lr){5-6} \cmidrule(lr){8-9} \cmidrule(lr){10-11}
& & \textbf{BP-to-EP} & \textbf{EP-to-BP} & \textbf{BP-to-EP} & \textbf{EP-to-BP} & & \textbf{BP-to-EP} & \textbf{EP-to-BP} & \textbf{BP-to-EP} & \textbf{EP-to-BP} \\
\midrule
Base protocol & 70.76 & 38.56 & 38.69 & \textbf{0.4748} & \textbf{0.4828} & 67.63 & 86.45 & 86.70 & 0.0978 & 0.0981 \\
Task-balanced batches & 69.74 & 38.56 & 38.55 & 0.4751 & 0.4859 & 68.52 & 86.04 & 86.45 & 0.0987 & 0.1001 \\
Added-token path & \textbf{72.93} & 38.36 & 38.67 & 0.4768 & 0.4864 & 66.66 & 86.23 & 86.60 & 0.0986 & 0.1006 \\
Restricted classification loss & 70.74 & \textbf{38.70} & \textbf{38.72} & 0.4750 & 0.4832 & \textbf{69.10} & \textbf{87.01} & \textbf{87.32} & \textbf{0.0941} & \textbf{0.0956} \\
\bottomrule
\end{tabular}
\end{table*}

No candidate consistently improves identification and rewriting across both
evaluation resources. The alternatives instead improve isolated metrics.
Task-balanced batches improve Golden Collection identification while reducing
the other reported scores. Added-token treatment improves FRMT identification
but not rewriting or Golden Collection identification. Restricted loss gives
the strongest Golden Collection values and a small FRMT BLEU increase, but it
does not improve FRMT identification. Since these improvements are confined to individual metrics while requiring
additional batching, tokenizer, or loss mechanisms, the base protocol is
retained as the simpler, broadly competitive configuration for the subsequent
experiments with the 4B-parameter model. The retained protocol therefore uses the
natural rewriting-to-identification mixture, the existing \texttt{PT}/\texttt{BR}
base-vocabulary representations, and normal decoder-vocabulary scoring at the
supervised identification label position, while rewriting uses loss over the
complete target sequence.

\subsection{Supervised Identification and Rewriting Results}
\label{sec:main-supervised-results}

With the supervised protocol selected by the 270M-parameter study, the
4B-parameter T5Gemma 2 model is adapted with rank-48 LoRA under both control
formulations. \Cref{tab:r48-classification-breakdown,tab:r48-translation-breakdown}
report identification and rewriting results after Stage B and Stage C. This
section discusses the supervised Stage B results, while
\Cref{sec:reward-continuation-results} examines the subsequent Stage C
continuation.

\begin{table*}[tbp]
\centering
\small
\renewcommand{\arraystretch}{1.12}
\setlength{\tabcolsep}{5pt}
\caption{Identification results for the 4B-parameter T5Gemma 2 model with rank-48 LoRA.}
\label{tab:r48-classification-breakdown}
\begin{tabular}{@{}clp{5.1cm}cccc@{}}
\toprule
\textbf{Stage} & \textbf{Control} & \textbf{Data condition} &
\multicolumn{2}{c}{\textbf{FRMT}} &
\multicolumn{2}{c}{\textbf{Golden Collection}} \\
\cmidrule(lr){4-5} \cmidrule(lr){6-7}
& & & \shortstack{\textbf{Accuracy}\\\textbf{(\%)}}
& \shortstack{\textbf{Macro-F1}\\\textbf{(\%)}}
& \shortstack{\textbf{Accuracy}\\\textbf{(\%)}}
& \shortstack{\textbf{Macro-F1}\\\textbf{(\%)}} \\
\midrule
B & Encoder & GPT-Wikipedia & 82.39 & 82.39 & \textbf{73.99} & \textbf{73.92} \\
B & Decoder & GPT-Wikipedia & 81.69 & 81.68 & 71.61 & 71.53 \\
B & Encoder & GPT-Wikipedia + FRMT & 86.09 & 86.06 & 71.43 & 71.19 \\
B & Decoder & GPT-Wikipedia + FRMT & \textbf{86.15} & \textbf{86.10} & 72.53 & 72.38 \\
\midrule
C & Encoder & GPT-Wikipedia & 82.33 & 82.33 & 73.63 & 73.50 \\
C & Decoder & GPT-Wikipedia & 81.52 & 81.31 & 72.34 & 71.81 \\
C & Encoder & GPT-Wikipedia + FRMT & 85.82 & 85.78 & 71.43 & 71.17 \\
C & Decoder & GPT-Wikipedia + FRMT & 85.65 & 85.54 & 72.34 & 72.02 \\
\bottomrule
\end{tabular}
\end{table*}

Stage B identification first shows how adding FRMT supervision changes the
two evaluation resources. On FRMT, macro-F1 rises from 82.39 to 86.06 under
encoder control and from 81.68 to 86.10 under decoder control. On the Golden
Collection, the effect is mixed: encoder-controlled macro-F1 decreases from
73.92 to 71.19, while decoder-controlled macro-F1 increases from 71.53 to
72.38. Adding FRMT therefore improves FRMT identification under both
formulations but does not produce a consistent Golden Collection gain, and
neither control formulation is consistently stronger across both evaluation
resources.

For external context, the encoder-controlled GPT-Wikipedia Stage B
configuration reaches 82.39 macro-F1 on FRMT, compared with the 77.25 F1
reported by \citet{Sousa_Almeida_Silvano_Cantante_Campos_Jorge_2025} for their
strongest BERTimbau-based identifier. Although both studies use FRMT for
evaluation, the comparison is contextual, as Sousa et al. use a dedicated
identification model with a different training setup and objective, while the
present configuration learns identification and rewriting in the same model.

\begin{table*}[tbp]
\centering
\small
\renewcommand{\arraystretch}{1.12}
\setlength{\tabcolsep}{3pt}
\caption{Directional rewriting results for the 4B-parameter T5Gemma 2 model after Stage B and Stage C with rank-48 LoRA. Lower TER is better.}
\label{tab:r48-translation-breakdown}
\begin{tabular}{@{}clp{4.4cm}cccc@{}}
\toprule
\textbf{Stage} & \textbf{Control} & \textbf{Data condition} &
\multicolumn{2}{c}{\textbf{FRMT}} &
\multicolumn{2}{c}{\textbf{Golden Collection}} \\
\cmidrule(lr){4-5} \cmidrule(lr){6-7}
& & & \textbf{BP-to-EP} & \textbf{EP-to-BP} & \textbf{BP-to-EP} & \textbf{EP-to-BP} \\
& & & \shortstack{\textbf{BLEU (0--100)}\\\textbf{/ TER}} & \shortstack{\textbf{BLEU (0--100)}\\\textbf{/ TER}} & \shortstack{\textbf{BLEU (0--100)}\\\textbf{/ TER}} & \shortstack{\textbf{BLEU (0--100)}\\\textbf{/ TER}} \\
\midrule
B & Encoder & GPT-Wikipedia & 37.87/0.4791 & 40.10/0.4714 & 86.08/0.0946 & \textbf{89.09}/\textbf{0.0823} \\
B & Encoder (rewriting-only) & GPT-Wikipedia & 37.73/0.4803 & 40.31/0.4715 & 85.66/0.0976 & 88.66/0.0847 \\
B & Decoder & GPT-Wikipedia & 40.45/0.4594 & 40.60/0.4673 & 87.75/0.0874 & 87.50/0.0916 \\
B & Encoder & GPT-Wikipedia + FRMT & 37.99/0.4770 & \textbf{45.43}/0.4273 & 65.14/0.2555 & 66.61/0.2535 \\
B & Encoder (rewriting-only) & GPT-Wikipedia + FRMT & 36.76/0.4852 & 45.40/\textbf{0.4251} & 56.10/0.3222 & 58.69/0.3103 \\
B & Decoder & GPT-Wikipedia + FRMT & \textbf{43.97}/\textbf{0.4338} & 44.74/0.4306 & 68.37/0.2336 & 67.83/0.2400 \\
\midrule
C & Encoder & GPT-Wikipedia & 37.89/0.4784 & 40.28/0.4696 & 86.22/0.0953 & 88.93/0.0837 \\
C & Decoder & GPT-Wikipedia & 40.34/0.4603 & 40.33/0.4698 & \textbf{88.44}/\textbf{0.0824} & 88.02/0.0878 \\
C & Encoder & GPT-Wikipedia + FRMT & 38.35/0.4756 & 43.47/0.4440 & 77.16/0.1601 & 79.95/0.1488 \\
C & Decoder & GPT-Wikipedia + FRMT & 42.72/0.4417 & 42.72/0.4466 & 81.52/0.1351 & 80.70/0.1402 \\
\bottomrule
\end{tabular}
\end{table*}

Across the Stage B comparisons, data composition produces larger changes in
rewriting than in identification. Adding FRMT changes FRMT rewriting differently across directions and control
formulations, ranging from almost no change in some comparisons to a clear
improvement in others. In contrast, Golden Collection BLEU drops by roughly 20
points across every direction and control comparison, with corresponding
increases in TER.

The rewriting-only baselines test whether identification supervision harms
rewriting. With GPT-Wikipedia, the multitask and rewriting-only configurations
remain very close: multitask BLEU is higher in three of the four
dataset-direction comparisons, while rewriting-only is higher only on FRMT
EP-to-BP, by 0.21 BLEU. With GPT-Wikipedia + FRMT, the multitask configuration
obtains higher BLEU in all four comparisons, including gains of 9.04 points in
Golden Collection BP-to-EP and 7.92 points in EP-to-BP. Across the eight
directional BLEU comparisons, rewriting-only training exceeds multitask
training only once and by a small margin, so removing identification
supervision does not improve rewriting overall.

Control placement shows a clearer directional tendency: decoder control
performs better in every BP-to-EP rewriting comparison, whereas EP-to-BP is
mixed and includes settings where encoder control performs better.

The most directly comparable prior rewriting result is BP-to-EP evaluation on
the Golden Collection. After rescaling the BLEU values reported by
\citet{sanches2024brazilianportugueseeuropeanportuguese} to the 0--100 scale,
their strongest trained system obtains approximately 84 BLEU and 0.10 TER,
while their published single-direction copy baseline reaches approximately 89
BLEU and 0.07 TER. The decoder-controlled GPT-Wikipedia configuration reaches
87.75 BLEU and 0.0874 TER after Stage B and 88.44 BLEU and 0.0824 TER after
Stage C. The comparison is contextual, as Sanches et al. train specifically for
BP-to-EP rewriting, while the present configuration is bidirectional and also
learns identification. \Cref{sec:appendix-external-model-comparison} provides a broader comparison with external pretrained and instruction-following systems. Across these systems, the best adapted T5Gemma 2 result is stronger in every reported identification and rewriting setting except FRMT BP-to-EP, where AMALIA SFT obtains 44.65 BLEU and 0.4289 TER compared with 43.97 BLEU and 0.4338 TER for the strongest adapted configuration.

After the Stage B comparisons, \Cref{tab:stage-a-ablation} isolates the
contribution of Stage A OpenSubtitles warm-up by comparing each configuration
with its corresponding no-Stage-A variant. The clearest benefit appears after
GPT-Wikipedia specialization under encoder control. On FRMT, Stage A adds 0.87
macro-F1 points and 4.40 BLEU points while reducing TER by 0.0573. On the
Golden Collection, the same configuration gains 4.01 macro-F1 points and 25.41
BLEU points while reducing TER by 0.2404. Under decoder control with
GPT-Wikipedia, Stage A has almost no rewriting effect on FRMT and lowers
macro-F1 by 1.28 points, but improves all three aggregate measures on the
Golden Collection. The GPT-Wikipedia + FRMT conditions show smaller or mixed
changes: encoder control improves FRMT but lowers Golden Collection rewriting
by 6.79 BLEU points, whereas decoder control produces small rewriting gains on
both resources. Stage A therefore provides its clearest aggregate benefit after
GPT-Wikipedia specialization, especially on the Golden Collection, rather than
improving every configuration uniformly.

\begin{table*}[tbp]
\centering
\small
\renewcommand{\arraystretch}{1.12}
\setlength{\tabcolsep}{5pt}
\caption{Stage A warm-up ablation under aggregate evaluation. Deltas are computed as the Stage A model minus the corresponding no-Stage-A model. Positive deltas indicate improvement for macro-F1 and BLEU, while negative deltas indicate improvement for TER.}
\label{tab:stage-a-ablation}
\begin{tabular}{@{}llp{3.2cm}ccc@{}}
\toprule
\textbf{Control} & \textbf{Data condition} & \textbf{Dataset}
& \shortstack{\(\Delta\) \textbf{Macro-F1}\\\textbf{(pp)}}
& \shortstack{\(\Delta\) \textbf{BLEU}\\\textbf{(points)}}
& \(\Delta\) \textbf{TER} \\
\midrule
Encoder & GPT-Wikipedia & FRMT & \textbf{+0.87} & \textbf{+4.40} & \textbf{-0.0573} \\
Encoder & GPT-Wikipedia & Golden Collection & \textbf{+4.01} & \textbf{+25.41} & \textbf{-0.2404} \\
Encoder & GPT-Wikipedia + FRMT & FRMT & +0.81 & +1.92 & -0.0176 \\
Encoder & GPT-Wikipedia + FRMT & Golden Collection & +1.62 & -6.79 & +0.0425 \\
Decoder & GPT-Wikipedia & FRMT & -1.28 & -0.01 & +0.0004 \\
Decoder & GPT-Wikipedia & Golden Collection & +2.67 & +1.42 & -0.0144 \\
Decoder & GPT-Wikipedia + FRMT & FRMT & -0.07 & +0.35 & -0.0031 \\
Decoder & GPT-Wikipedia + FRMT & Golden Collection & +2.19 & +2.00 & -0.0217 \\
\bottomrule
\end{tabular}
\end{table*}

\subsection{Reward-Based Continuation}
\label{sec:reward-continuation-results}

\Cref{tab:r48-classification-breakdown,tab:r48-translation-breakdown} also
report how the Stage B configurations change after Stage C. Although
continuation uses rewriting-only data, identification remains close to the
corresponding Stage B values. Macro-F1 changes are small across both resources
and control formulations, with increases and decreases rather than a consistent
improvement.

The rewriting effect depends strongly on the Stage B supervision condition.
After GPT-Wikipedia specialization, Stage C changes rewriting only modestly,
with BLEU differences below one point across all four dataset-direction
comparisons. The larger effect follows GPT-Wikipedia + FRMT specialization.
Golden Collection BLEU recovers by approximately 12--13 points in both
directions under both control formulations, while FRMT rewriting changes remain
smaller and differ across directions and formulations. Stage C therefore has
its clearest effect as a recovery of Golden Collection rewriting after
GPT-Wikipedia + FRMT specialization rather than as a uniform improvement across
evaluation resources.

These changes characterize the complete Stage C procedure. Since continuation
changes both the optimization procedure and the rewriting subset used after
Stage B, the reported comparison does not separate the effect of reward-based
optimization from the effect of the continuation data.

\subsection{Qualitative Analysis}
\label{sec:qualitative-analysis}

The qualitative analysis examines how source-reference similarity and model
copying behavior relate to the rewriting differences observed above. The
Golden Collection contains a high proportion of exact or near-copy
references. Of its 1,000 bidirectional rows, 454 are exact source-reference
matches, 566 have copy TER at most 0.05, and 636 have copy TER at most 0.08.
Consistent with this composition, the unchanged-source baseline reaches 89.17
BLEU and 0.0755 TER on the Golden Collection. The corresponding FRMT baseline is substantially weaker, at 38.75 BLEU and
0.4819 TER, consistent with larger source-reference differences in that
resource. This contrast helps explain why unnecessary rewriting can be costly on the
Golden Collection while FRMT provides more room to improve over unchanged-source
copying.

These reference profiles also clarify why useful local adaptations do not
necessarily raise corpus BLEU above the Golden Collection copy baseline. A
required edit can improve only a limited span, whereas an unnecessary
paraphrase can remove many matching higher-order \(n\)-grams. Aggregate BLEU
can therefore favor the unchanged source even when local edits are appropriate,
but this does not make copying universally desirable.

Encoder-controlled Stage B models make this contrast between conservative and
more extensive rewriting visible. On the 454 exact-copy rows, the model trained with
GPT-Wikipedia data preserves the source in 74.7\% of cases, compared with
18.7\% for the model trained with the GPT-Wikipedia + FRMT data mix. On these
rows, preserving the source is consistent with the reference. The model trained with GPT-Wikipedia data also remains more conservative on
the 546 non-exact rows, copying the source in 33.7\% of cases compared with
3.1\% for the model trained with the GPT-Wikipedia + FRMT data mix. On these
rows, copying can indicate under-editing when the reference contains a required
variety-specific change, although a source-reference difference can also
reflect paraphrastic or editorial changes that are not required for EP/BP
adaptation. Supervision from the GPT-Wikipedia + FRMT data mix sharply reduces
copying, which can help realize required changes but also creates more
opportunities for unnecessary reformulation.
\Cref{tab:golden-qualitative-overtranslation} illustrates how the two Stage B supervision conditions translate into different rewriting behavior.

In the first example, the reference replaces \textit{procura} with
\textit{demanda}. The
model trained with GPT-Wikipedia data leaves the source unchanged despite the required lexical
edit, whereas the model trained with the GPT-Wikipedia + FRMT data mix makes the change and reformulates the
sentence more broadly. In the second example, the reference is identical to the source.
The model trained with GPT-Wikipedia data preserves it, while the model trained
with the GPT-Wikipedia + FRMT data mix
unnecessarily reformulates this otherwise unchanged sentence.

\begin{table*}[tbp]
\centering
\small
\renewcommand{\arraystretch}{1.12}
\setlength{\tabcolsep}{3pt}
\caption{Representative Golden Collection outputs from encoder-controlled Stage B models trained with GPT-Wikipedia and GPT-Wikipedia + FRMT. Boldface marks model-output spans illustrating under-editing or unnecessary reformulation.}
\label{tab:golden-qualitative-overtranslation}
\begin{tabular}{@{}p{0.205\textwidth}p{0.205\textwidth}p{0.225\textwidth}p{0.225\textwidth}@{}}
\toprule
\textbf{Source} & \textbf{Reference} & \textbf{GPT-Wikipedia} & \textbf{GPT-Wikipedia + FRMT} \\
\midrule
\textit{``A procura aumenta de forma dr\'astica no mundo todo'', afirmou o especialista em recursos h\'idricos Jos\'e Galizia Tundisi, do Instituto Internacional de Ecologia, em S\~ao Carlos (SP).} &
\textit{``A demanda aumenta de forma dr\'astica no mundo todo'', afirmou o especialista em recursos h\'idricos Jos\'e Galizia Tundisi, do Instituto Internacional de Ecologia, em S\~ao Carlos (SP).} &
\textit{``A \textbf{procura} aumenta de forma dr\'astica no mundo todo'', afirmou o especialista em recursos h\'idricos Jos\'e Galizia Tundisi, do Instituto Internacional de Ecologia, em S\~ao Carlos (SP).} &
\textit{``A demanda \textbf{est\'a aumentando drasticamente em todo o mundo}'', \textbf{disse} o especialista em recursos h\'idricos Jos\'e Galizia Tundisi, do Instituto Internacional de Ecologia em S\~ao Carlos (SP).} \\

\textit{Da\'i por que est\'a absolutamente correto no conte\'udo, como na forma, o parecer do ilustre Relator Bernardo Cabral.} &
\textit{Da\'i por que est\'a absolutamente correto no conte\'udo, como na forma, o parecer do ilustre Relator Bernardo Cabral.} &
\textit{Da\'i por que est\'a absolutamente correto no conte\'udo, como na forma, o parecer do ilustre Relator Bernardo Cabral.} &
\textit{\textbf{Por isso, o parecer do ilustre Relator, Bernardo Cabral, \'e absolutamente correto tanto no conte\'udo como na forma.}} \\
\bottomrule
\end{tabular}
\end{table*}

A separate issue concerns the references used for evaluation. Some source-reference pairs contain differences that extend beyond localized EP/BP conversion, so a model can produce an appropriate variety-specific rewrite without reproducing every transformation present in the reference. \Cref{tab:evaluation-reference-information-differences} gives representative cases from the Golden Collection and FRMT. In the Golden Collection example, the reference removes the final clause describing the returned gesture, while in the FRMT example it removes the feet conversions while preserving the corresponding measurements in metres. Other inspected FRMT references contain lexical substitutions such as \textit{marcante} to \textit{importante}, \textit{ganhou} to \textit{recebeu}, \textit{reside} to \textit{vive}, \textit{usar} to \textit{utilizar}, and \textit{música} to \textit{canção}, which are valid alternatives but are not consistently required for variety adaptation. These examples show that reference construction can include paraphrastic or editorial differences that complicate reference-based evaluation in addition to the supervision-dependent model behavior illustrated in \Cref{tab:golden-qualitative-overtranslation}.

\begin{table*}[tbp]
\centering
\small
\renewcommand{\arraystretch}{1.12}
\setlength{\tabcolsep}{3pt}
\caption{Examples where source and reference sentences differ in expressed information beyond localized variety conversion. Boldface marks source spans whose information is omitted from the reference.}
\label{tab:evaluation-reference-information-differences}
\begin{tabular}{@{}p{0.13\textwidth}p{0.09\textwidth}p{0.34\textwidth}p{0.34\textwidth}@{}}
\toprule
\textbf{Dataset} & \textbf{Direction} & \textbf{Source} & \textbf{Reference} \\
\midrule
Golden Collection &
BP-to-EP &
\textit{- Oww voc\^e \'e um amor Obrigado viu!- \textbf{Falei e o mesmo retribuiu com um abra\c{c}o.}} &
\textit{- Oww, \'es ador\'avel. Obrigado!} \\

FRMT &
BP-to-EP &
\textit{O Aut\'odromo Internacional Orlando Moura tem uma pista de 3.433 metros \textbf{(11.263 p\'es)}, e o Kart\'odromo Ayrton Senna uma pista de 930 metros \textbf{(3.051 p\'es)}.} &
\textit{O Aut\'odromo Internacional Orlando Moura tem uma pista de 3433 metros e o Kart\'odromo Ayrton Senna tem uma pista de 930 metros.} \\
\bottomrule
\end{tabular}
\end{table*}

\section{Discussion and Limitations}
\label{sec:discussion-limitations}

The supervised and qualitative results show that aligned EP/BP resources are not interchangeable simply by virtue of containing valid parallel sentences. Supervision composition produces substantially larger changes in rewriting than in identification, and \Cref{tab:golden-qualitative-overtranslation} shows that the GPT-Wikipedia and GPT-Wikipedia + FRMT conditions induce markedly different rewriting behavior. GPT-Wikipedia was constructed around localized variety adaptation, while FRMT provides a broader region-aware translation resource. Under the evaluated conditions, the usefulness of aligned supervision therefore depends on how closely the transformations represented in the training pairs match the intended rewriting behavior, not only on the amount of parallel data available.

Control placement has a smaller and less uniform effect than supervision composition. Decoder control is consistently stronger in the reported BP-to-EP rewriting comparisons, whereas EP-to-BP remains mixed and includes settings where encoder control performs better. Identification likewise provides no universal ordering between the two formulations across the evaluation resources. The experiments establish these resource- and direction-dependent differences but do not determine why placing variety information on the encoder or decoder side changes their relative performance, so control placement remains an empirical design choice rather than a uniformly preferable formulation.

The staged comparisons reinforce the dependence on the behavior learned before an additional adaptation step. Stage A provides its clearest aggregate benefit after GPT-Wikipedia specialization, especially on the Golden Collection, while the other evaluated conditions show smaller or mixed changes. Stage C changes rewriting only modestly after GPT-Wikipedia specialization but substantially restores Golden Collection performance after GPT-Wikipedia + FRMT specialization, with smaller and mixed changes on FRMT. These patterns indicate that an additional training stage modifies a rewriting policy established by the preceding supervision rather than providing a uniform gain across starting conditions and evaluation resources.

These conclusions are also constrained by the construction of the held-out evaluation references and by the metrics used to compare outputs with them. As the source-reference examples in \Cref{sec:qualitative-analysis} show, a valid evaluation pair can contain paraphrastic, editorial, or information-level differences beyond the localized variety changes required by the task, while the Golden Collection also contains many cases that require little or no rewriting. BLEU and TER remain useful surface-oriented measures, but a single reference cannot reliably distinguish every required variety change, appropriate copying, acceptable alternative phrasing, and unnecessary reformulation. The same reference dependence also enters Stage C, where sentence-level BLEU and TER against the reference contribute directly to the reward, so broader-than-required reference changes can also influence the continuation signal.

The experimental scope introduces additional limitations. All reported experiments use the T5Gemma 2 encoder-decoder family, so the effects of supervision composition, control placement, and staged adaptation have not been established across multiple model families. The 270M-parameter full-fine-tuning study and the experiments with the 4B-parameter model serve different purposes and use different adaptation methods, so they do not constitute a controlled comparison of scale or fine-tuning strategy. Stage C also changes the continuation subset and optimization procedure together, which prevents their individual contributions from being separated in the reported comparison. These limitations define the scope within which the quantitative and qualitative findings should be interpreted.

\section{Conclusion}
\label{sec:conclusion}

This work examined Portuguese variety identification and bidirectional EP/BP rewriting within a shared T5Gemma 2 text-to-text setup, with particular attention to how supervision composition, control placement, and staged adaptation shape the two tasks. The experiments show that one adapted model can support both identification and rewriting, while rewriting behavior is especially sensitive to the transformations represented in the aligned supervision. GPT-Wikipedia, which was constructed around localized variety adaptation, produces conservative rewriting that aligns closely with the Golden Collection, whereas adding FRMT changes the learned rewriting behavior substantially and produces different effects across resources and directions.

Control placement and staged adaptation refine this picture without providing a single universally preferable configuration. Adding FRMT consistently improves FRMT identification under both control formulations, while its effect on Golden Collection identification is mixed. For rewriting, decoder control is consistently stronger for BP-to-EP, whereas EP-to-BP remains mixed across the evaluated conditions, and removing identification supervision does not improve rewriting overall. Stage A provides its clearest benefit in selected GPT-Wikipedia conditions, while Stage C changes the GPT-Wikipedia condition only modestly and has its largest effect after GPT-Wikipedia + FRMT specialization, where it restores a substantial part of the Golden Collection rewriting performance. The adapted configurations also numerically exceed the closest published task-specific results used for comparison, namely Sousa et al.~\citeyearpar{Sousa_Almeida_Silvano_Cantante_Campos_Jorge_2025} on FRMT identification and the strongest trained Sanches et al.~\citeyearpar{sanches2024brazilianportugueseeuropeanportuguese} system on Golden Collection BP-to-EP rewriting, although these comparisons involve dedicated setups that differ from the shared bidirectional setting used here. Across the additional external systems evaluated in \Cref{sec:appendix-external-model-comparison}, the best adapted result is stronger in every reported identification and rewriting setting except FRMT BP-to-EP, where AMALIA SFT is slightly stronger.

The broader implication is that effective EP/BP rewriting depends both on supervision whose transformations reflect the intended variety-adaptation task and on evaluation resources and metrics that can distinguish required edits from appropriate preservation and acceptable alternative phrasing. Future work should therefore prioritize more tightly verified variety-focused supervision together with human evaluation, targeted linguistic checks, multiple acceptable references, or validated learned evaluators that better capture these distinctions. Further experiments should also separate continuation-data selection from reward-based optimization, extend the comparison to additional model families and domains, and include a current frontier proprietary model under the same identification and bidirectional rewriting protocol.

\bibliographystyle{ACM-Reference-Format}
\bibliography{sample-base}

\clearpage
\appendix

\section{Supplementary Data, Comparison, and Diagnostics}
\label{sec:appendix-supplementary-details}

This appendix provides supplementary information on resource construction, the
secondary no-equal diagnostic, a descriptive external-model comparison, and
reward-design diagnostics.
\Cref{sec:appendix-data-construction} describes the preparation of the training
resources, with particular attention to GPT-Wikipedia, which was constructed
for this work. \Cref{sec:appendix-no-equal-paper} reports the no-equal
diagnostic, \Cref{sec:appendix-external-model-comparison} reports the
external-model comparison, and \Cref{sec:appendix-rl-diagnostics} covers
reward-design sensitivity.

\subsection{Data Construction and Processing}
\label{sec:appendix-data-construction}

Resources from different origins are converted into a common aligned-pair
representation containing an EP sentence, its BP counterpart, provenance
metadata, and a split label. This representation supports both rewriting
directions through the same data interface. Non-equal pairs also supply
input-variety identification examples, whereas equal pairs remain useful for
rewriting and copying supervision. The common representation allows
resources prepared through different procedures to enter the same downstream
task-construction interface without conflating their supervisory roles.

GPT-Wikipedia is constructed from Portuguese Wikipedia article dumps used to
supply topics and contexts rather than copied aligned examples. Informative
leading sentences are retained from sufficiently developed, sentence-like
paragraphs, while short, list-like, or punctuation-poor material is excluded
from the inspiration batches. OpenAI GPT-4o, accessed through the IAedu API,
then generates new EP/BP pairs. The Portuguese generation instructions
require equivalent meaning, keep the two variants close, favor limited lexical
or morphosyntactic contrasts, and discourage changes not needed for variety
adaptation or extensive paraphrase. Wikipedia therefore supplies topical
material, while the aligned pair content is automatically generated under
constraints for localized EP/BP contrasts.

Generated pairs are audited at the level of changed token spans. Automatic
checks normalize superficial differences, identify unexpected substitutions or
reorderings, and assign reason categories to the resulting mappings. Recurring
transformation mappings are then aggregated for human review. The reviewer
approves mappings that represent valid EP/BP contrasts and rejects changes that
are unnecessary for variety adaptation. These decisions are propagated to
matching examples, with contextual exceptions handled separately. Approved
contrasts retain the appropriate EP and BP realizations, while rejected
differences are normalized to shared wording. This review reduces supervision
for unrelated paraphrase while preserving genuine variety contrasts and equal
pairs that indicate when rewriting is unnecessary.

OpenSubtitles is processed for broad Stage A exposure rather than as a tightly
controlled rewriting resource. Exact duplicates are removed, and
subtitle-specific formatting is stripped before fragmented sentences are
repaired or discarded. Adjacent rows are merged when continuation evidence in
either aligned sentence, or the pipeline's under-segmentation heuristic,
indicates that a sentence boundary was introduced by subtitle segmentation. A
punctuation-and-capitalization check prevents merges when both rows instead
appear to contain complete sentence boundaries. Contextual word-alignment
coverage and relative-length criteria then filter weak alignments, with
stronger evidence required for longer segments while retaining valid short
subtitle fragments. These operations preserve the resource's scale and
contextual coverage while reducing evident segmentation and alignment noise.

FRMT is used primarily as held-out evaluation data. Its training split is
introduced only in the GPT-Wikipedia + FRMT Stage B condition. It is
converted into the common pair format with its regional bucket metadata
retained, although the reported evaluation aggregates over the complete FRMT
split. The corresponding Stage C condition continues
from that Stage B checkpoint but uses a much smaller, more conservative
rewriting-only subset.
Consequently, the Stage C subset preserves the broader data condition without
being identical to the Stage B mixture.

\subsection{Secondary No-Equal Ablation}
\label{sec:appendix-no-equal-paper}

The decoder-controlled no-equal diagnostic removes equal source-reference
rewriting examples from the supervised stages while retaining the same Stage C
continuation data. Stage C therefore differs through its supervised starting
checkpoint rather than through a separate continuation-data variant. Aggregate
evaluation shows only small changes on FRMT but lower Golden Collection
rewriting performance without equal examples, supporting their retention in the
supervised data.

\subsection{External Model Comparison}
\label{sec:appendix-external-model-comparison}

Four pretrained external systems provide descriptive context for the
identification and rewriting results: AMALIA-9B-0626-SFT and
AMALIA-9B-0626-DPO \citep{simplicio2026amalia}, Qwen3-4B
\citep{yang2025qwen3technicalreport}, and FLAN-T5-XL \citep{3722577.3722647}. These models are
evaluated without task-specific fine-tuning on the EP/BP training resources
used for the T5Gemma 2 configurations. Rewriting uses direction-specific
Portuguese prompts and the same BLEU and TER implementations as the internal
evaluation, with BP-to-EP and EP-to-BP reported separately. For identification,
AMALIA uses free generation followed by label parsing, while Qwen3-4B and
FLAN-T5-XL use natural-label sequence scoring.
The evaluated checkpoints are
\path{amalia-llm/AMALIA-9B-0626-SFT},
\path{amalia-llm/AMALIA-9B-0626-DPO}, \path{Qwen/Qwen3-4B}, and
\path{google/flan-t5-xl}.
The comparison is descriptive, as the systems differ in architecture, scale,
training history, prompting, and inference procedure.

\begin{table*}[tbp]
  \centering
  \small
  \renewcommand{\arraystretch}{1.08}
  \setlength{\tabcolsep}{6pt}

  \caption{External-model results on FRMT and the Golden Collection.}

  \label{tab:appendix-external-model-comparison}

  \begin{tabular}{@{}llccc@{}}
    \toprule
    \textbf{Model}
    & \textbf{Dataset}
    & \shortstack{\textbf{Macro-F1}\\\textbf{(\%)}}
    & \textbf{BP-to-EP}
    & \textbf{EP-to-BP} \\
    & & & \shortstack{\textbf{BLEU (0--100)}\\\textbf{/ TER}} & \shortstack{\textbf{BLEU (0--100)}\\\textbf{/ TER}} \\
    \midrule

    AMALIA SFT
    & FRMT
    & 45.37
    & \textbf{44.65}/\textbf{0.4289}
    & 40.22/0.4682 \\

    & Golden Collection
    & 43.99
    & 68.44/0.2347
    & 74.66/0.1920 \\

    \addlinespace[0.2em]

    AMALIA DPO
    & FRMT
    & \textbf{59.72}
    & 43.82/0.4494
    & 35.16/0.6063 \\

    & Golden Collection
    & \textbf{56.55}
    & 58.08/0.3239
    & 57.18/0.3773 \\

    \addlinespace[0.2em]

    Qwen3-4B
    & FRMT
    & 52.56
    & 39.44/0.4709
    & \textbf{40.73}/\textbf{0.4654} \\

    & Golden Collection
    & 49.96
    & \textbf{78.33}/\textbf{0.1540}
    & \textbf{78.26}/\textbf{0.1590} \\

    \addlinespace[0.2em]

    FLAN-T5-XL
    & FRMT
    & 47.87
    & 32.22/0.5263
    & 32.17/0.5373 \\

    & Golden Collection
    & 40.99
    & 67.06/0.2209
    & 66.31/0.2292 \\

    \bottomrule
  \end{tabular}
\end{table*}

AMALIA DPO obtains the highest external identification Macro-F1 on both
evaluation resources, with 56.55 on the Golden Collection and 59.72 on FRMT.
These values remain below the strongest adapted results of 73.92 and 86.10,
respectively, in \Cref{tab:r48-classification-breakdown}.

On the Golden Collection, Qwen3-4B is the strongest external rewriting system
in both directions. It obtains 78.33 BLEU and .1540 TER for BP-to-EP, compared
with 88.44 BLEU and .0824 TER for the Stage C decoder-controlled
GPT-Wikipedia configuration. For EP-to-BP, Qwen3-4B obtains 78.26 BLEU and
.1590 TER, compared with 89.09 BLEU and .0823 TER for the Stage B
encoder-controlled GPT-Wikipedia configuration. The adapted T5Gemma 2
model is stronger in both Golden Collection directions in
\Cref{tab:r48-translation-breakdown}.

On FRMT, the strongest external system changes by direction. For BP-to-EP,
AMALIA SFT obtains 44.65 BLEU and .4289 TER, slightly outperforming the Stage B
decoder-controlled GPT-Wikipedia + FRMT configuration at 43.97 BLEU and .4338 TER in
both metrics. For EP-to-BP, Qwen3-4B obtains 40.73 BLEU and .4654 TER, while the
Stage B encoder-controlled GPT-Wikipedia + FRMT configuration reaches 45.43 BLEU and
.4273 TER. The strongest adapted T5Gemma 2 configuration therefore outperforms
Qwen3-4B in this direction. The comparison is direction-dependent and does not
support a single ordering between the external and adapted models.

\subsection{Reward-Design Diagnostics}
\label{sec:appendix-rl-diagnostics}

The fully fine-tuned 270M reward-design diagnostics compare four configurations
independently. These are the base BLEU/TER reward, augmentation with a
classifier reward, augmentation with token conservatism, and TER-diverse
candidate grouping. Both encoder-controlled and decoder-controlled
formulations are evaluated. Each diagnostic begins independently from the same
corresponding supervised Stage B checkpoint and changes only the named reward
component or candidate-group construction relative to the base diagnostic.

The base reward scores each candidate using BLEU and TER relative to the
reference rewrite and the unchanged-source baseline. For decoder-controlled
diagnostics, an incorrect first \texttt{PT}/\texttt{BR} control label sets the
candidate reward to zero. This first-label gate is a binary check of the
generated control token. The classifier variant instead adds a continuous
signal obtained by evaluating the generated sentence in identification mode
and using the model probability assigned to the requested target variety. It
therefore measures how strongly the generated sentence itself is recognized as
the target variety. The token-conservatism variant adds token-level reward for
reference-supported changes and partial credit for reference-supported
copying. TER-diverse grouping changes candidate selection rather than adding a
reward component. It selects candidates spanning different TER values relative
to the reference while retaining the base reward.

\begin{table}[H]
\centering
\small
\renewcommand{\arraystretch}{1.08}
\setlength{\tabcolsep}{4pt}
\caption{Decoder-controlled 270M isolated Stage C ablations.}
\label{tab:appendix-rl-diagnostics-decoder-paper}
\begin{tabular}{@{}lccc@{}}
\toprule
\textbf{Configuration}
& \shortstack{\textbf{Macro-F1}\\\textbf{(\%)}}
& \shortstack{\textbf{BLEU}\\\textbf{(0--100)}}
& \textbf{TER} \\
\midrule
\multicolumn{4}{@{}l}{\textbf{FRMT}} \\
Supervised Stage B start & \textbf{70.76} & \textbf{39.13} & \textbf{0.4788} \\
\midrule
Base reward & 68.76 & 39.11 & 0.4800 \\
Classifier reward & 69.54 & 39.04 & 0.4806 \\
Token conservatism & 69.49 & 39.07 & 0.4804 \\
TER-diverse grouping & 68.24 & 39.12 & 0.4799 \\
\addlinespace[0.35em]
\multicolumn{4}{@{}l}{\textbf{Golden Collection}} \\
Supervised Stage B start & \textbf{67.63} & 86.57 & 0.0980 \\
\midrule
Base reward & 67.15 & 86.88 & 0.0962 \\
Classifier reward & 67.57 & 86.96 & 0.0969 \\
Token conservatism & 67.43 & 86.58 & 0.1007 \\
TER-diverse grouping & 66.70 & \textbf{87.19} & \textbf{0.0949} \\
\bottomrule
\end{tabular}
\end{table}

\begin{table}[H]
\centering
\small
\renewcommand{\arraystretch}{1.08}
\setlength{\tabcolsep}{4pt}
\caption{Encoder-controlled 270M isolated Stage C ablations.}
\label{tab:appendix-rl-diagnostics-encoder-paper}
\begin{tabular}{@{}lccc@{}}
\toprule
\textbf{Configuration}
& \shortstack{\textbf{Macro-F1}\\\textbf{(\%)}}
& \shortstack{\textbf{BLEU}\\\textbf{(0--100)}}
& \textbf{TER} \\
\midrule
\multicolumn{4}{@{}l}{\textbf{FRMT}} \\
Supervised Stage B start & 71.73 & \textbf{27.06} & \textbf{0.6204} \\
\midrule
Base reward & 71.90 & 13.62 & 0.7525 \\
Classifier reward & 71.81 & 15.67 & 0.7319 \\
Token conservatism & \textbf{72.28} & 16.27 & 0.7268 \\
TER-diverse grouping & 72.22 & 13.18 & 0.7571 \\
\addlinespace[0.35em]
\multicolumn{4}{@{}l}{\textbf{Golden Collection}} \\
Supervised Stage B start & 65.43 & \textbf{65.53} & \textbf{0.2996} \\
\midrule
Base reward & 65.77 & 38.21 & 0.5154 \\
Classifier reward & 66.29 & 44.25 & 0.4704 \\
Token conservatism & \textbf{66.34} & 42.79 & 0.4808 \\
TER-diverse grouping & 66.14 & 38.96 & 0.5118 \\
\bottomrule
\end{tabular}
\end{table}
\FloatBarrier

Under decoder control, FRMT rewriting remains extremely close to the supervised
Stage B starting point across the four diagnostics. TER-diverse grouping gives
the strongest Golden Collection rewriting among the diagnostic configurations,
with BLEU 87.19 and TER 0.0949 compared with 86.57 and 0.0980 at the Stage B
start. Its Golden Collection Macro-F1 is lower, at 66.70 compared with 67.63.
No diagnostic improves identification and rewriting consistently across both
evaluation resources.

Under encoder control, all four configurations substantially reduce rewriting
quality relative to the supervised Stage B start. On FRMT, the starting BLEU is
27.06, while the best configuration reaches 16.27 with token conservatism. On
the Golden Collection, the starting BLEU is 65.53, while the best configuration
reaches 44.25 with classifier reward. Identification changes are comparatively
small, and token conservatism gives the highest diagnostic Macro-F1 on both
resources, at 72.28 on FRMT and 66.34 on the Golden Collection.

The two formulations begin from different supervised Stage B starting points.
Their changes should therefore be interpreted relative to their own baselines
rather than by comparing their absolute values directly.

\end{document}
\endinput
%%
%% End of file `sample-acmtog.tex'.
